\section{DESI-KL Forecast}
\label{sec:desi}

In \citep{xu2024KL} we simulate a potential future DESI-KL survey; the main results of this work will be summarized in this section.

A DESI-KL project would require placing multiple optical fibers along a disk galaxy's major and minor axes to constrain its velocity field. This concept has already been successfully implemented for the DESI Peculiar Velocity Survey (DESI-PV) \citep{DESIPV23}, which will be talked about in Sec. \ref{sec:pv}. We used the DESI Legacy Imaging Survey \citep{LS_overview} as the basis for our galaxy sample selection. The Legacy Survey covers $14,000$ deg$^{2}$ spanning approximately in declination $-18\degree< \delta < +84\degree$. The Legacy Survey Data Release 9\footnote{\url{https://www.legacysurvey.org/dr9/}}~\citep[LS DR9 from hereon][]{SDH+21} contains three optical bands in $(g,\,r,\,z)$ and four infrared bands from WISE. The optical bands images are composed by three survey projects: the Beijing-Arizona Sky Survey~\citep[BASS,][]{BASS}, the DECam Legacy Survey~\citep[DECaLS,][]{LS_overview}, and the Mayall $z$-band Legacy Survey~\citep[MzLS,][]{LS_overview}. We use the tractor sweep catalog ~\citep{LHM16,L20} for the photometry, shape, photo-$z$, and morphology measurements. The Bright Galaxy Survey (BGS) \citep{BGS_selection1, BGS_selection2} is another DESI survey containing central fiber measurements for already selected galaxies, which we can also use to constrain our samples.

For our galaxy sample selection, we imposed criteria with varying strictness on its brightness and morphology, as described below. For selecting bright galaxies, we have:
\begin{itemize}
    \item Criterion 1.a LS $r\leq22$: This is the most relaxed selection criterion among the three. We pose similar selection criteria as the one in~\citep{BGS_selection2} but relax the $r$-band magnitude cut to 22. The motivation is to include more galaxies to study the trade-off between number density v.s. SNR.  
    \item Criterion 1.b BGS Any: This is more stringent than 1.a and we include all galaxies that are either selected as BGS Bright or BGS Faint. The benefit is that their redshift and emission line properties are already measured by BGS, which eliminates many of the uncertainties in target selection using the photometric catalog. Also, we only need to take offset fiber exposures for the BGS Any sample. 
    \item Criterion 1.c BGS Bright: This criterion selects galaxies from the BGS Bright sample. This ensures high SNR but reduces the number density.
\end{itemize}
Then, for selecting disk galaxies, we have:
\begin{itemize}
    \item Criterion 2.a: The tractor morphology is classified as inclined exponential, round exponential, or Sersic profile and with Sersic index $\leq 2$, and the above-atmosphere hlr $\geq 1\arcsec$.
    \item Criterion 2.b: Same as 2.a but with a more conservative cut on the above-atmosphere hlr $\geq 1.\arcsec5$.
\end{itemize}
Groups of samples are created by combining one type of filter from Criteria 1 and one from Criteria 2, which are listed in Table \ref{tab:sample_def}. The abundance of galaxy samples from each selection criterion is shown in Fig. \ref{fig:zdist} in the form of its number density as a function of redshift.

For each set of samples, we then performed a simulated cosmological likelihood analysis for $\sigma_8$, and compared the constraints to LSST weak lensing (LSST-WL) mock analyses. The results for Samples 1 and 3 are shown in Fig. \ref{fig:sigma8}, where we see comparable or better constraints from the simulated DESI-KL analysis. Given that LSST and DESI will be observing mostly different areas of the sky, we have shown that DESI-KL would be a great complement to LSST-WL measuring cosmic shear over the entire extragalactic sky.

\begin{table}
    \centering
    \begin{tabular}{ccccc}
    \toprule
        Sample & \shortstack{Selection\\Criteria} & $n_\mathrm{gal}$ (arcmin$^{-2}$) & $z_0$ & $\alpha$\\
        \tableline 
        1 & 1.a \& 2.a & 0.527 & 0.0894 & 0.8826 \\
        2 & 1.a \& 2.b & 0.173 & 0.0259 & 0.6835 \\
        3 & 1.b \& 2.a & 0.180 & 0.1027 & 1.1616 \\
        4 & 1.b \& 2.b & 0.097 & 0.0569 & 0.9701 \\
        5 & 1.c \& 2.a & 0.120 & 0.0777 & 1.1031 \\
        6 & 1.c \& 2.b & 0.083 & 0.0504 & 0.9537 \\
        \tableline
    \end{tabular}
    \caption{Summary of KL sample target selection criteria and the resulting sample number density and redshift distribution. We define six samples with their target selection criteria shown in the second column and the surface number density in the third column. We show the best-fitting parameters of $n(z)$ distribution in Eq.~(\ref{eqn:nz_param}) in columns four and five.}
    \label{tab:sample_def}
\end{table}